\chapter{Introducción}
Esta plantilla sigue el estilo común que emplean los estudiantes de la \textit{Universidad de El Salvador} para sus trabajos de graduación: portada, aprobación de junta/directiva, agradecimientos, índices (figuras, tablas, algoritmos y códigos fuente), capítulos, apéndices y bibliografía.

\section{Estructura del proyecto}
La estructura sugerida de carpetas y archivos es:
\begin{verbatim}
.
|-- DEIthesis.cls
|-- appendices/
|   |-- appendix1.tex
|   `-- code_listings.tex
|-- chapters/
|   |-- 01_introduction.tex
|   `-- 02_conclusion.tex
|-- code/
|-- frontmatter/
|   |-- abstract.tex
|   `-- thanks.tex
|-- main.tex
|-- references.bib
|-- res/           % logos institucionales
`-- assets/        % imágenes/figuras del documento
\end{verbatim}

\section{Figuras, tablas, algoritmos y código}
En la \cref{fig:ejemplo-figura} se ilustra una figura de ejemplo. La \cref{tab:ejemplo-tabla} muestra una tabla. El \cref{alg:ejemplo} ilustra un algoritmo. El \cref{lst:ejemplo-codigo} muestra un fragmento de código.

\begin{figure}[H]
  \centering
  % Coloca aquí assets como PNG/JPG/PDF; si no existe, se dibuja un recuadro.
  \fbox{\rule{0pt}{4cm}\rule{6cm}{0pt}} % marcador de posición
  \caption{Figura de ejemplo}\label{fig:ejemplo-figura}
\end{figure}

\begin{table}[H]
  \centering
  \caption{Tabla de ejemplo}\label{tab:ejemplo-tabla}
  \begin{tabular}{lcc}
    \toprule
    Métrica & Valor 1 & Valor 2 \\
    \midrule
    A & 10 & 20 \\
    B & 30 & 40 \\
    \bottomrule
  \end{tabular}
\end{table}

\begin{algorithm}[H]
  \DontPrintSemicolon
  \caption{Algoritmo de ejemplo}\label{alg:ejemplo}
  \KwIn{Entrada $x$}
  \KwOut{Salida $y$}
  $y \leftarrow 0$\;
  \For{$i \leftarrow 1$ \KwTo $n$}{
    $y \leftarrow y + f(x_i)$\;
  }
  \Return{$y$}\;
\end{algorithm}

\begin{lstlisting}[language=Python,caption={Código de ejemplo en Python},label={lst:ejemplo-codigo}]
def suma(a, b):
    return a + b

print(suma(2, 3))
\end{lstlisting}

\section{Citas y bibliografía}
Ejemplo de cita \parencite{lamport94,knuth1984texbook}. El estilo activo es \texttt{apa} con backend \texttt{biber}.

\section{Acrónimos}
Ejemplo de acrónimo: \gls{ppg}.
Otro acrónimo: \gls{ecg}. Y otro más: \gls{latex}.